\chapter[Referencial Teórico]{Referencial Teórico}
\label{chap: fundamentacao}

Este capítulo apresenta o embasamento teórico do trabalho. A exposição parte da contextualização do mercado de jogos digitais e do fenômeno dos esports (\autoref{sec: mercado}), avança para o conceito de equilíbrio competitivo (\autoref{sec: equilibrio}) e para a noção de arquétipos em jogos de luta (\autoref{sec: arquetipos}), introduz os fundamentos dos Algoritmos Genéticos (\autoref{sec: ag}) e da otimização multi-objetivo via NSGA-II (\autoref{sec: nsga2}), e fecha com a revisão dos trabalhos relacionados que aplicam algoritmos evolutivos ao balanceamento de jogos (\autoref{sec: trabalhos_relacionados}).

\section{Jogos digitais competitivos e esports}
\label{sec: mercado}

A indústria de jogos digitais consolidou-se nas últimas duas décadas como um dos segmentos mais expressivos da economia do entretenimento. Segundo o relatório global da \citeonline{newzoo2025}, a receita mundial do setor atingiu US\$~188{,}8 bilhões em 2025, com projeção de US\$~206{,}5 bilhões em 2028, e uma base global estimada em 3{,}58 bilhões de jogadores --- aproximadamente sessenta por cento da população com acesso à internet. O contingente é distribuído de forma desigual entre plataformas, com participação predominante de jogos para dispositivos móveis, mas com retomada acentuada dos consoles e crescimento estável do segmento de computadores pessoais. No cenário nacional, a 13\textsuperscript{a} edição da \citeonline{pgb2026} indica que aproximadamente 75\% da população brasileira declara consumir jogos digitais como atividade regular, e que a América Latina figura como a quarta maior região do mundo em número de jogadores. Os dados são relatórios de indústria, e não fontes peer-reviewed; comparecem aqui por ausência de equivalentes acadêmicos com a mesma abrangência quantitativa, mas com a devida ressalva quanto à sua origem.

Sobre essa base econômica articula-se o fenômeno dos \textit{esports} --- competições estruturadas de jogos digitais com regras formalizadas, premiações, equipes profissionais e públicos próprios. \citeonline{hamari2017esports} propõem uma definição operacional que tornou-se referência no campo: esports são uma forma de esporte cujos aspectos primários são mediados por sistemas eletrônicos, tanto na entrada dos jogadores quanto na saída do sistema de jogo, sendo realizados sob interfaces humano-computador. A definição enfatiza o caráter mediado da prática --- diferenciando-a do esporte tradicional sem reduzi-la a mero entretenimento --- e fornece base conceitual para o tratamento acadêmico do fenômeno. \citeonline{reitman2020esports}, em revisão sistemática que cobre o intervalo entre 2002 e 2018, mapeiam o campo de pesquisa em esports como atividade distribuída por sete disciplinas distintas --- administração, ciências do esporte, ciências cognitivas, informática, direito, estudos de mídia e sociologia ---, indicando a consolidação dos esports como objeto multidisciplinar de investigação.

A relevância dos esports para o problema tratado neste trabalho é direta. Em jogos cuja proposta inclui dimensão competitiva organizada, o equilíbrio entre opções de jogo --- personagens, classes, decks --- deixa de ser apenas requisito de \textit{design} e passa a ser condição para a viabilidade da própria competição. Um elenco de personagens em que apenas um subconjunto seja competitivamente viável reduz drasticamente a profundidade estratégica do jogo, frustra investimentos de jogadores em personagens preteridos e compromete a credibilidade do sistema competitivo \cite{adams2014fundamentals}. O custo dessa exigência, todavia, é elevado: o balanceamento manual de elencos com dezenas de personagens demanda iterações sucessivas de desenvolvimento e teste, normalmente sustentadas por jogadores especializados, e frequentemente requer correções pós-lançamento à medida que a comunidade competitiva identifica configurações dominantes não previstas durante o desenvolvimento. A automação parcial desse processo, por meio de técnicas de busca e otimização, constitui motivação prática significativa para a presente investigação.

\section{Equilíbrio competitivo em jogos}
\label{sec: equilibrio}

O conceito de equilíbrio competitivo --- frequentemente referido pelo termo \textit{balance} na literatura --- não possui definição única e universalmente aceita. \citeonline{adams2014fundamentals} caracteriza-o, em sentido amplo, como a condição em que nenhum recurso, estratégia ou opção disponível ao jogador domina sistematicamente as alternativas, de modo que escolhas distintas permaneçam competitivamente viáveis. \citeonline{elias2012characteristics} expandem a discussão ao listar o equilíbrio entre as propriedades estruturais que organizam a análise comparativa de jogos, distinguindo equilíbrio entre jogadores --- que assegura igualdade de oportunidade entre partes em condições simétricas --- de equilíbrio interno ao espaço de opções --- que sustenta diversidade estratégica sem dominância.

Uma distinção conceitual particularmente relevante para este trabalho é a que separa equilíbrio entendido como \textbf{simetria estatística} de equilíbrio entendido como \textbf{estrutura cíclica de vantagens}. No primeiro caso, exige-se que cada opção apresente desempenho médio aproximadamente equivalente contra o conjunto das demais. No segundo, admite-se que opções individuais possam vencer ou perder confrontos específicos --- possivelmente de forma sistemática ---, contanto que o conjunto, considerado coletivamente, não permita que uma opção domine todas as outras. A estrutura cíclica é, em muitos gêneros, propriedade desejável: ela introduz não-transitividade ao espaço de opções, sustenta a importância da escolha contextual e gera variedade estratégica que a simetria estatística pura tende a apagar.

A formalização desse princípio em jogos de luta foi proposta por \citeonline{gingold2006rps}, em demonstração que articula o jogo \textit{Pedra-Papel-Tesoura} como caso-base mínimo de equilíbrio não-transitivo e, por construção, deriva variantes em tempo real cuja estrutura preserva a propriedade cíclica original. O argumento estabelece que jogos de luta competitivos podem ser entendidos como extensões em tempo contínuo de \textit{Pedra-Papel-Tesoura} ampliado, em que as ações disponíveis e seus tempos de execução diferenciam-se, mas a estrutura subjacente de vantagens cíclicas permanece como fundamento do desenho competitivo. Trata-se de um resultado conceitual com consequências práticas: o equilíbrio buscado em jogos de luta não é, em geral, a equalização das taxas de vitória individuais, e sim a preservação de uma estrutura cíclica em que cada personagem encontre vantagens e desvantagens recíprocas distribuídas pelo elenco.

A consequência metodológica para este trabalho é direta. Equilibrar um elenco de personagens não significa fazer todas as taxas de vitória convergirem para cinquenta por cento contra todas as demais; significa garantir que nenhuma taxa de vitória apresente desvio acentuado em relação a esse valor de referência, sem que isso requeira a destruição da estrutura cíclica de vantagens recíprocas que caracteriza o gênero. Essa distinção orienta a formulação da função de aptidão adotada no \autoref{chap: metodologia}, em que a dominância competitiva entre personagens é medida sem que o ciclo de vantagens seja codificado como restrição.

\section{Arquétipos em jogos de luta}
\label{sec: arquetipos}

Em jogos de luta, o termo \textbf{arquétipo} designa um agrupamento conceitual de personagens que compartilham um modo característico de jogar, derivado de uma combinação reconhecível de atributos físicos --- como velocidade, alcance, dano e resistência --- e de uma estratégia preferencial de uso desses atributos durante o combate. A noção é amplamente empregada pela comunidade competitiva de jogos de luta para classificar personagens, orientar escolhas de equipamento e descrever vantagens entre confrontos, mas não possui, até onde foi possível verificar, definição taxonômica consolidada na literatura acadêmica peer-reviewed. As fontes que delimitam arquétipos individuais com algum grau de detalhe são, predominantemente, materiais produzidos por jogadores profissionais, sítios eletrônicos especializados e fóruns de discussão competitiva --- não tratáveis como referências formais sob critérios estritos de revisão por pares.

Diante desse cenário, este trabalho adota uma postura explícita: o recorte de cinco arquétipos descrito nesta seção é uma \textbf{construção operacional} do autor, derivada da convenção de uso corrente na comunidade competitiva, e ancorada teoricamente apenas no princípio cíclico estabelecido por \citeonline{gingold2006rps}. Cada arquétipo é operacionalizado neste trabalho como um vetor de atributos numéricos e pesos comportamentais, especificados no \autoref{chap: metodologia}, e nenhum dos cinco rótulos adotados deve ser entendido como categoria taxonômica formalmente reconhecida na literatura. A escolha de cinco --- e não três, número mínimo para sustentar um ciclo --- é deliberada: cinco arquétipos constituem o menor elenco em que cada personagem vence dois confrontos e perde dois, configurando uma estrutura cíclica com profundidade combinatória suficiente para que a evolução dos atributos por Algoritmo Genético produza resultados não-triviais, sem aproximar excessivamente o problema da formulação canônica de \textit{Pedra-Papel-Tesoura}.

\subsection{Os cinco arquétipos adotados}

Os cinco arquétipos modelados neste trabalho são caracterizados a seguir, em descrição funcional --- não operacional. A operacionalização numérica completa, com bounds dos atributos e valores canônicos iniciais, é apresentada no \autoref{chap: metodologia}.

\begin{itemize}
    \item \textit{Rushdown} --- personagem orientado à ofensiva contínua e à pressão de curta distância. Caracteriza-se por alta velocidade e baixo tempo de recuperação entre ataques, com alcance reduzido. Sua estratégia preferencial é fechar a distância rapidamente, manter o oponente sob pressão e impedir que ele estabeleça ações que dependam de preparação.

    \item \textit{Zoner} --- personagem orientado ao controle de espaço a longa distância. Caracteriza-se por alcance elevado, dano moderado a longo prazo e capacidade de empurrar o oponente para fora da zona de punição. Sua estratégia preferencial é manter o adversário afastado, dificultando a aproximação.

    \item \textit{Grappler} --- personagem orientado ao dano explosivo a curta distância. Caracteriza-se por dano máximo por golpe, resistência elevada e mobilidade reduzida. Sua estratégia preferencial é absorver dano enquanto se aproxima e converter a aproximação em punição decisiva.

    \item \textit{Combo Master} --- personagem orientado à conversão de acertos isolados em sequências de dano. Caracteriza-se por capacidade elevada de imobilizar temporariamente o oponente após um golpe bem-sucedido, prolongando a janela em que pode continuar atacando sem retaliação. Sua estratégia preferencial é equilibrar exposição e oportunidade até obter o acerto inicial.

    \item \textit{Turtle} --- personagem orientado à defesa atritiva. Caracteriza-se por resistência elevada, capacidade defensiva acentuada e tempo de recuperação prolongado entre ataques. Sua estratégia preferencial é absorver a pressão adversária, aguardar erros do oponente e capitalizar sobre eles.
\end{itemize}

\subsection{O ciclo de vantagens canônico}

Os cinco arquétipos descritos articulam-se em uma estrutura cíclica de vantagens em que cada arquétipo vence dois confrontos e perde dois, conforme o \autoref{quad: ciclo_canonico}. As vantagens descritas refletem a operacionalização adotada neste trabalho e foram justificadas por argumentos funcionais derivados da convenção de uso na comunidade competitiva; não pretendem reproduzir, exata e universalmente, todos os matchups observados em jogos de luta comerciais, e algumas das relações aqui codificadas são objeto de leituras divergentes na comunidade. A discussão pormenorizada dessas divergências, bem como o tratamento metodológico do ciclo como hipótese a ser avaliada \textit{a posteriori}, são desenvolvidos no \autoref{chap: metodologia}.

\begin{quadro}[ht]
    \caption{Ciclo canônico de vantagens entre os cinco arquétipos adotados.}
    \label{quad: ciclo_canonico}
    \centering

    \begin{tabularx}{\textwidth}{>{\raggedright\arraybackslash}p{2.4cm} >{\raggedright\arraybackslash}p{3.3cm} >{\raggedright\arraybackslash}p{3.3cm} >{\raggedright\arraybackslash}X}
        \hline
        \bfseries Vencedor & \bfseries Perde para & \bfseries Vence & \bfseries Justificativa funcional \\
        \hline
        \textit{Rushdown}     & \textit{Grappler}, \textit{Turtle}     & \textit{Zoner}, \textit{Combo Master}   & A pressão contínua impede que o adversário estabeleça ações que dependam de preparação. \\
        \textit{Zoner}        & \textit{Rushdown}, \textit{Combo Master} & \textit{Grappler}, \textit{Turtle}      & O controle do espaço a longa distância mantém o adversário fora da zona de punição. \\
        \textit{Grappler}     & \textit{Zoner}, \textit{Combo Master}  & \textit{Rushdown}, \textit{Turtle}      & O dano explosivo a curta distância pune o recuo do agressor e neutraliza a postura defensiva. \\
        \textit{Combo Master} & \textit{Rushdown}, \textit{Turtle}     & \textit{Zoner}, \textit{Grappler}       & A conversão de um acerto em sequência prolongada é decisiva contra alvos lentos ou frágeis. \\
        \textit{Turtle}       & \textit{Zoner}, \textit{Grappler}      & \textit{Rushdown}, \textit{Combo Master} & A defesa atritiva absorve a pressão e neutraliza sequências que dependem do acerto inicial. \\
        \hline
    \end{tabularx}

    \source
\end{quadro}

É essencial sublinhar que o ciclo canônico apresentado no \autoref{quad: ciclo_canonico} não é codificado como restrição no processo evolutivo descrito nos capítulos subsequentes. Ele cumpre dois papéis distintos. Em primeiro lugar, fornece a semente inicial da população, sob a forma dos valores canônicos dos atributos associados a cada arquétipo. Em segundo lugar, serve como métrica de avaliação \textit{a posteriori}, permitindo verificar em que medida a estrutura cíclica de vantagens é preservada, modificada ou destruída ao longo da evolução. Codificar o ciclo diretamente na função de aptidão tornaria a investigação circular: o sistema preservaria a estrutura porque foi pago para preservá-la, e não porque ela emerge das mecânicas de combate e da pressão simultânea por equilíbrio e por preservação de identidade.

\section{Algoritmos Genéticos}
\label{sec: ag}

Algoritmos Genéticos (AGs) constituem uma classe de algoritmos de busca e otimização inspirados nos princípios da evolução biológica, formalizados por \citeonline{holland1975adaptation} e consolidados como técnica computacional madura por \citeonline{goldberg1989genetic}. A ideia central é manter uma população de soluções candidatas a um problema --- denominadas \textbf{indivíduos} --- e fazê-la evoluir, ao longo de sucessivas \textbf{gerações}, mediante operadores inspirados em fenômenos biológicos: seleção dos indivíduos mais aptos, recombinação de seu material genético e introdução de variação aleatória. O processo continua até que algum critério de parada seja atingido, momento em que o melhor indivíduo encontrado é retornado como solução aproximada ao problema.

\subsection{Representação do indivíduo}

Cada indivíduo é representado por um \textbf{cromossomo}, codificação em geral vetorial das variáveis de decisão do problema. Cada posição do cromossomo é um \textbf{gene}, e o conjunto de valores admissíveis para cada gene define o \textbf{espaço de busca} explorado pelo algoritmo. A representação pode ser binária, inteira, real ou estruturada, conforme a natureza do problema; representações reais são usuais em problemas de otimização paramétrica contínua, e representações estruturadas tornam-se necessárias quando o problema envolve composições heterogêneas de variáveis, como ocorre quando atributos de natureza distinta convivem em um mesmo indivíduo \cite{eiben2015introduction,mitchell1998introduction}.

\subsection{Função de aptidão}

A \textbf{função de aptidão} (\textit{fitness function}) atribui a cada indivíduo um valor numérico que reflete a qualidade da solução por ele representada, conforme os critérios do problema. Em problemas de minimização, valores menores indicam soluções melhores; em maximização, o inverso. A função de aptidão é o único elo formal entre o algoritmo e o domínio do problema --- toda informação utilizada na evolução é mediada por ela ---, razão pela qual sua formulação é frequentemente a decisão de \textit{design} mais sensível em aplicações de AGs. Quando o problema envolve múltiplos critérios não diretamente comparáveis, a função de aptidão pode ser construída como soma ponderada das componentes individuais, ou como vetor de objetivos avaliados em paralelo, abordagem que conduz à otimização multi-objetivo discutida na \autoref{sec: nsga2}.

\subsection{Operadores evolutivos}

A evolução da população é conduzida por três operadores principais.

A \textbf{seleção} escolhe, entre os indivíduos da geração corrente, aqueles que contribuirão para a geração seguinte, com probabilidade tipicamente proporcional à aptidão. Diversos mecanismos de seleção foram propostos na literatura, entre os quais a seleção por roleta, a seleção por torneio e mecanismos baseados em \textit{ranking} \cite{eiben2015introduction}. A seleção por torneio, em particular, sorteia subconjuntos de indivíduos do tamanho especificado por um parâmetro e seleciona o mais apto de cada subconjunto, oferecendo controle simples sobre a pressão seletiva.

O \textbf{cruzamento} (\textit{crossover}) combina o material genético de dois indivíduos selecionados para produzir descendentes. Esquemas de cruzamento variam conforme a representação adotada e o grau de coerência interna que se deseja preservar entre genes correlacionados. Operadores de cruzamento por ponto único ou por blocos de genes são frequentes em representações cuja semântica admite que subconjuntos de genes sejam internamente coerentes e devam ser herdados em conjunto \cite{mitchell1998introduction}.

A \textbf{mutação} introduz pequenas perturbações aleatórias nos genes dos descendentes, com probabilidade controlada por um parâmetro denominado \textbf{taxa de mutação}. Em representações reais, a perturbação é tipicamente gaussiana, centrada no valor corrente do gene, com desvio-padrão fixado como fração da amplitude admissível para aquele gene. A mutação é o operador responsável por introduzir diversidade genética que o cruzamento sozinho não consegue gerar, e sua calibração equilibra exploração do espaço de busca e refinamento de soluções já promissoras --- o trade-off entre exploração e explotação descrito por \citeonline{eiben1998exploration}.

\subsection{Elitismo e convergência}

O \textbf{elitismo} é um mecanismo complementar pelo qual um subconjunto dos melhores indivíduos da geração corrente é transmitido sem alterações para a geração seguinte, impedindo que soluções de alta qualidade já encontradas sejam perdidas por efeito dos operadores estocásticos. A fração da população preservada por elitismo é parâmetro de \textit{design}, tipicamente pequena em relação ao tamanho total.

A execução do algoritmo prossegue até que um critério de parada seja atendido. Critérios usuais incluem o atingimento de um número máximo de gerações, a estagnação da melhor aptidão por número de gerações pré-determinado e o cumprimento de uma condição específica do problema --- como, no caso deste trabalho, a obtenção de equilíbrio competitivo abaixo de um limiar pré-especificado em todos os confrontos. A combinação dessas condições é discutida no \autoref{chap: metodologia}.

\section{Otimização multi-objetivo e NSGA-II}
\label{sec: nsga2}

Problemas de otimização envolvendo um único critério são chamados de \textbf{escalares} ou \textbf{mono-objetivos}; problemas envolvendo dois ou mais critérios são denominados \textbf{multi-objetivos}. Quando os critérios são parcialmente conflitantes --- isto é, quando melhorar um deles tende a degradar outro ---, a noção tradicional de solução ótima única deixa de ser aplicável, e o problema passa a admitir um conjunto de soluções qualitativamente distintas que representam diferentes compromissos entre os objetivos \cite{deb2001multiobjective}.

\subsection{Escalarização e suas limitações}

Uma abordagem comum para tratar problemas multi-objetivos consiste em reduzi-los a problemas escalares por meio da \textbf{escalarização}: as componentes individuais são combinadas em uma única função de aptidão, em geral por soma ponderada, e o problema resultante é resolvido por técnicas mono-objetivo padronizadas. A escalarização é simples de implementar e permite o uso de algoritmos consolidados, mas apresenta limitações conhecidas. Em primeiro lugar, requer que o projetista escolha previamente os pesos das componentes, decisão que frequentemente carece de justificativa principiada e que determina \textit{a priori} o ponto da curva de compromisso que o algoritmo buscará. Em segundo lugar, varreduras sistemáticas no espaço de pesos podem deixar de identificar regiões da curva de compromisso quando esta apresenta concavidades. Em terceiro lugar, e talvez mais importante, a escalarização colapsa em um único número uma informação que é intrinsecamente multi-dimensional, ocultando do projetista a natureza do compromisso entre os objetivos \cite{deb2001multiobjective}.

\subsection{Dominância de Pareto e fronteira de Pareto}

A formalização multi-objetivo propriamente dita opera diretamente sobre o vetor de objetivos. Diz-se que uma solução $a$ \textbf{domina} uma solução $b$ no sentido de Pareto quando $a$ é ao menos tão boa quanto $b$ em todos os objetivos e estritamente melhor em pelo menos um deles. Soluções não dominadas por nenhuma outra do conjunto compõem a \textbf{fronteira de Pareto} do problema, conjunto que representa o lugar geométrico dos compromissos não dominados entre os objetivos. Cada ponto da fronteira corresponde a uma solução qualitativamente distinta, e a curva agregada descreve o trade-off entre os objetivos sem necessidade de definir pesos relativos \cite{deb2001multiobjective}.

\subsection{NSGA-II}

O \textit{Non-dominated Sorting Genetic Algorithm II} (NSGA-II), proposto por \citeonline{deb2002nsga}, é o algoritmo multi-objetivo evolutivo mais amplamente adotado na literatura. Sua arquitetura combina três mecanismos principais: a ordenação não-dominada rápida, que classifica a população em níveis sucessivos de não-dominância; a distância de aglomeração, que estima a densidade local de soluções na fronteira e promove diversidade ao penalizar regiões super-povoadas; e a seleção elitista por torneio binário com base nos dois critérios anteriores. O resultado é uma aproximação progressiva da fronteira de Pareto, com soluções distribuídas ao longo dela com diversidade preservada explicitamente.

A vantagem do NSGA-II para problemas com tensão entre critérios é que ele expõe a fronteira de Pareto de forma explícita, em vez de colapsar o compromisso entre objetivos em uma única solução escalar. No contexto deste trabalho, em que o equilíbrio competitivo e a preservação de identidade arquetípica constituem objetivos parcialmente conflitantes, a aplicação do NSGA-II permite caracterizar a forma do trade-off entre essas duas dimensões e selecionar representantes específicos da fronteira --- como o ponto de menor dominância, o ponto de menor desvio ou o ponto de máxima curvatura --- conforme o critério desejado, sem comprometimento prévio com uma ponderação escalar arbitrária.

\section{Trabalhos relacionados: balanceamento de jogos via algoritmos evolutivos}
\label{sec: trabalhos_relacionados}

A aplicação de algoritmos evolutivos ao balanceamento de jogos é tema investigado há aproximadamente duas décadas, com trabalhos que cobrem domínios variados. Esta seção revisa os trabalhos mais próximos do problema tratado neste estudo e situa a contribuição da presente investigação em relação a eles.

\citeonline{hom2007automatic} aplicaram Algoritmos Genéticos ao projeto automatizado de jogos de tabuleiro equilibrados, codificando as regras de jogos abstratos como cromossomos e avaliando indivíduos por meio de simulações com agentes computacionais. O trabalho demonstrou a viabilidade da abordagem evolutiva para problemas de \textit{design} de jogos --- estabelecendo um precedente metodológico relevante --- mas operou sobre um domínio em que as regras do jogo são parte do espaço de busca, e não sobre o balanceamento dos parâmetros de personagens dentro de regras fixas.

\citeonline{browne2010evolutionary} expandiram a abordagem evolutiva no contexto do \textit{evolutionary game design}, descrevendo a geração automática de jogos completos por evolução de suas regras, com avaliação por meio de critérios estruturais e de simulação. O escopo é mais amplo que o do balanceamento de elenco propriamente dito, e o objeto evoluído é a estrutura do jogo, não o conjunto de seus personagens.

\citeonline{preuss2012balanced} aplicaram busca evolutiva à determinação de planos balanceados em um cenário competitivo específico, demonstrando o uso de coevolução para identificar configurações em que a competição entre os agentes preserva diversidade estratégica. O foco recai sobre os planos de jogo, não sobre a especificação dos personagens em si.

O trabalho mais próximo do problema aqui tratado é o de \citeonline{chen2014mmorpg}, que aplicaram programação coevolutiva ao balanceamento das funções de crescimento de atributos por raça em \textit{Massively Multiplayer Online Role-Playing Games} de combate por turnos. As similaridades com o presente trabalho são significativas: o problema é descrito como o balanceamento entre classes de personagens em \textit{Player versus Player}, e a abordagem é evolutiva, com avaliação por simulação de combates. Há, contudo, três diferenças metodológicas importantes que sustentam o caráter complementar do presente trabalho. Em primeiro lugar, o domínio é distinto: jogos de \textit{role-playing} massivos por turnos não envolvem o aspecto posicional contínuo nem as janelas de ação sub-segundo características dos jogos de luta em tempo real, o que impõe requisitos diferentes sobre o modelo de simulação. Em segundo lugar, a formulação multi-objetivo via NSGA-II não é empregada; o problema é tratado em formulação coevolutiva escalar. Em terceiro lugar, e talvez mais relevante, a preservação de identidade entre as raças não é objeto formal da função de aptidão --- o balanceamento é o único critério explicitamente otimizado, e a manutenção da diferenciação entre classes é tratada como propriedade desejável apenas implicitamente.

\citeonline{togelius2011searchbased} oferecem uma sistematização mais ampla, em forma de \textit{survey}, das aplicações de busca à geração procedural de conteúdo em jogos. O trabalho situa o balanceamento como uma das categorias dentro de um campo mais amplo de busca aplicada a jogos, fornecendo o referencial conceitual em que se inserem os estudos específicos discutidos acima.

A combinação dos elementos abordados pelos trabalhos revisados --- balanceamento por algoritmos evolutivos, simulação de combate como meio de avaliação, e tratamento de classes de personagens --- constitui a base sobre a qual a presente investigação se ergue. A lacuna que o trabalho preenche reside na intersecção de três requisitos não atendidos simultaneamente por nenhum dos estudos anteriores: o domínio dos jogos de luta em tempo real, a formulação multi-objetivo com exposição explícita do trade-off por meio da fronteira de Pareto, e a preservação de identidade arquetípica como objetivo formalmente medido e otimizado, ortogonal ao equilíbrio.